\documentclass{article}
\usepackage{booktabs}
\usepackage{tabularx}
\usepackage[margin=1in]{geometry}
\begin{document}

\begin{table}[tbp] \centering
\newcolumntype{R}{>{\raggedleft\arraybackslash}X}
\newcolumntype{L}{>{\raggedright\arraybackslash}X}
\newcolumntype{C}{>{\centering\arraybackslash}X}

\caption{Federal viticulture subsidy line item, 1900-1920 (HSSO I.33a + I.33b, 1000 CHF)}
\label{tab:viti_subsidy}
\begin{tabularx}{\linewidth}{@{}lCC@{}}

\toprule
&{(1)}&{(2)} \tabularnewline \midrule
{Year}&{Viticulture subsidy (1000 CHF)}&{Era / event marker} \tabularnewline
\midrule \addlinespace[\belowrulesep]
1900&(blank)&no line item
1901&(blank)&no line item
1902&(blank)&no line item
1903&(blank)&no line item
1904&(blank)&no line item
1905&(blank)&no line item
1906&(blank)&no line item
1907&(blank)&no line item
1908&133&FIRST APPEARS (ban year)
1909&237&1909 implementing ordinance
1910&277&
1911&115&
1912&130&
1913&251&
1914&243&WWI budget
1915&43&WWI budget
1916&90&WWI budget
1917&54&WWI budget
1918&59&WWI budget
1919&67&
1920&190&
\bottomrule \addlinespace[\belowrulesep]

\end{tabularx}
\\ \parbox{\linewidth}{\footnotesize Notes: Federal viticulture subsidy line item from HSSO I.33a column C (1866-1915) and I.33b column H (continuation). The line is BLANK from the start of the federal agricultural subsidy series in 1866 through 1907 and first appears in 1908 with 133K CHF -- the same calendar year as the absinthe-ban referendum (5 July 1908, vote \\#68). The 1909 row reflects the implementing ordinance of the 1906 Lebensmittelgesetz that established federal thujone limits. Two interpretations of the 1908 emergence remain possible without Brugger 1968 disambiguation: (a) a genuinely new federal program emerged in 1908 (full-spectrum-capture interpretation: positive transfers + competitor elimination in the same year); or (b) a pre-existing program newly broken out from a previously-aggregated category (categorical-reclassification interpretation). The post-1908 trajectory (237K 1909, 277K 1910) shows the line item growing in scale, and the post-vote continuation in I.33b is the foundation for the substrate-substitution industrial policy documented in stage (iv) of the policy-economy arc. The post-1908 wine-index surge (1910 = 107 vs. 1905 \~ 50) coincides with the ban-imposed elimination of cheap-absinthe substitute demand. This pattern is consistent with regulatory-capture-induced demand redirection toward wine, distinct from the natural supply-recovery from phylloxera replanting (which occurred 1895-1905). The two sources of wine-market improvement are temporally separable and substantively distinct. WWI budget reallocation explains the 1914-1918 dip. Source: HSSO I.33a + I.33b, citing Brugger 1968. See progress\\_2026-05-11\\_2215\\_viti1908.md for the full discovery context.}
\end{table}
\end{document}
